\documentclass{article}
\usepackage{mathtools} 
\usepackage{fontspec}
\usepackage[UTF8]{ctex}
\usepackage{amsthm}
\usepackage{mdframed}
\usepackage{xcolor}
\usepackage{amssymb}
\usepackage{amsmath}


% 定义新的带灰色背景的说明环境 zremark
\newmdtheoremenv[
  backgroundcolor=gray!10,
  % 边框与背景一致，边框线会消失
  linecolor=gray!10
]{zremark}{说明}

% 通用矩阵命令: \flexmatrix{矩阵名}{元素符号}{行数}{列数}
\newcommand{\flexmatrix}[4]{
  \[
  #1 = \begin{pmatrix}
    #2_{11}     & #2_{12}     & \cdots & #2_{1#4}   \\
    #2_{21}     & #2_{22}     & \cdots & #2_{2#4}   \\
    \vdots      & \vdots      & \ddots & \vdots     \\
    #2_{#31}    & #2_{#32}    & \cdots & #2_{#3#4}
  \end{pmatrix}
  \]
}

% 简化版命令（默认矩阵名为A，元素符号为a）: \quickmatrix{行数}{列数}
\newcommand{\quickmatrix}[2]{\flexmatrix{A}{a}{#1}{#2}}

\begin{document}
\title{注释 17.2}
\author{张志聪}
\maketitle

\begin{zremark}
  先抛出来一个问题：

  复合函数求导，为什么可以按照，树状图的路径走？
\end{zremark}
\textbf{证明：}

直接看定理17.5本身，不容易看出来原因，要从其证明过程入手。

改造一下$x$，设$x = \varphi(s)$
于是在树状图中$x$将没有$t$的路径，这是因为$x$的微分表示将变为
\begin{align}
  \Delta = \frac{\partial x}{\partial x} \Delta s + \alpha_1 \Delta s
\end{align}
缺少了$\frac{\partial x}{\partial t}$部分。


\end{document}